\newpage
\noindent\textbf{APÊNDICE A -- REGISTRO DO AMBIENTE DE INTELIGÊNCIA ARTIFICIAL (ANTIGRAVITY)}
\vspace{0.5cm}

\noindent Este apêndice registra a instalação, a configuração e o uso do ambiente de Inteligência Artificial baseado em agentes \textbf{Google Antigravity}, empregado no apoio à elaboração, estruturação metodológica, busca bibliográfica e revisão técnica desta etapa da pesquisa de mestrado.

\vspace{0.3cm}
\noindent\textbf{FERRAMENTA INSTALADA}

\noindent Utilizou-se o assistente de desenvolvimento baseado em agentes \textbf{Antigravity}, proporcionando recursos avançados de refatoração em \LaTeX, auditoria textual, organização modular do projeto e conformidade com as diretrizes do Programa de Pós-Graduação em Engenharia de Controle e Automação (ProPECAut/Ifes).

\vspace{0.3cm}
\noindent\textbf{CONFIGURAÇÃO DO AMBIENTE}

\noindent A configuração concentrou-se no espaço de trabalho local do projeto de mestrado:
\begin{enumerate}[label=\alph*)]
    \item Arquitetura modular com pacote de estilo customizado (\texttt{config.sty});
    \item Separação estruturada dos elementos textuais (capa, resumo, objetivos e notas metodológicas);
    \item Adequação das entregas entre os marcos de qualificação e defesa da dissertação;
    \item Organização de base bibliográfica especializada em sistemas de suprimento de energia espacial e eletrônica de potência.
\end{enumerate}

\vspace{0.3cm}
\noindent\textbf{USO E APLICAÇÃO NESTE TRABALHO}

\noindent O assistente Antigravity atuou em regime de cooperação e escrita conjunta (\textit{pair writing}) com o pesquisador nas seguintes tarefas:
\begin{itemize}[leftmargin=*]
    \item \textbf{Alinhamento Temático e Metodológico:} Estruturação das frentes de pesquisa voltadas ao Subsistema de Suprimento de Energia (\textit{EPS}), eficiência de conversores estáticos e algoritmos de MPPT para plataformas de pico e nanossatélites;
    \item \textbf{Definição de Metas e Entregas:} Formulação da matriz de objetivos específicos divididos entre qualificação e defesa, vinculados às respectivas entregas e produtos tecnológicos;
    \item \textbf{Pesquisa Bibliográfica Inicial:} Mapeamento do estado da arte em arquiteturas de microrredes espaciais, modelos fotovoltaicos sob gradientes de LEO e técnicas de gerenciamento dinâmico de carga.
\end{itemize}